\documentclass[11pt]{article}

\usepackage[T1]{fontenc}
\usepackage[margin=1.1in]{geometry}
\usepackage{amsmath,amssymb,amsthm}
\usepackage{enumitem}
\usepackage{booktabs}
\usepackage[colorlinks=true,linkcolor=blue,citecolor=blue,urlcolor=blue]{hyperref}

\newtheorem{theorem}{Theorem}[section]
\newtheorem{proposition}[theorem]{Proposition}
\newtheorem{lemma}[theorem]{Lemma}
\newtheorem{corollary}[theorem]{Corollary}
\newtheorem{definition}[theorem]{Definition}
\newtheorem{conjecture}[theorem]{Conjecture}
\newtheorem{question}[theorem]{Question}
\theoremstyle{remark}
\newtheorem{remark}[theorem]{Remark}
\newtheorem{example}[theorem]{Example}

\newcommand{\R}{\mathbb{R}}
\newcommand{\X}{X}
\newcommand{\Xhat}{\widehat{\X}}
\newcommand{\Mx}{M_{\X}}
\newcommand{\Mg}{M^{g}_{\X}}
\newcommand{\Rx}{\mathcal{R}_{\X}}
\newcommand{\Rg}{\mathcal{R}^{g}_{\X}}
\newcommand{\Hb}[1]{\mathcal{H}_{#1}}
\newcommand{\dist}{\operatorname{dist}}
\newcommand{\dX}{d_{\X}}
\newcommand{\conv}{\operatorname{conv}}
\newcommand{\hconv}{\operatorname{hconv}}
\newcommand{\inj}{\operatorname{inj}}
\newcommand{\sect}{\operatorname{sec}}
\newcommand{\sca}[2]{\left\langle #1,\, #2\right\rangle}
\newcommand{\E}{\mathbb{E}}
\newcommand{\HH}{\mathbb{H}}
\newcommand{\Sn}{\mathbb{S}^{n}}
\newcommand{\bdry}{\partial_{\infty}}

\title{Medial-axis reconstruction on complete Riemannian
manifolds:\\ a self-contained theory of capture, drift, and the
unconditional characterisation\thanks{This paper develops, from a
single external starting point --- Bia\l{}o\.zyt's Euclidean
Theorem~6 \cite{Bia} --- the complete generalisation of medial-axis
reconstruction to arbitrary complete Riemannian manifolds. It
consolidates and re-derives, in one self-contained account, results
the author previously distributed across a sequence of companion
notes on Hadamard, hyperbolic, spherical, and mixed geometries; only
Bia\l{}o\.zyt's theorem is assumed as prior mathematics, and it is
used only to identify the \emph{faithful} certificate in the flat
case. Drafted with the assistance of Claude (Anthropic); all
arguments were re-derived and checked by hand, and the two mechanics
of the capture theorem --- the rate-one identity off the axis and the
margin at first approach --- were verified numerically on
$\mathbb{S}^{2}$ against a non-circular convex contact curve whose
focal region forces the foot to slide (60 random starts: growth rate
$1$ to five digits, margin deficit $<3\cdot10^{-6}$; file
\texttt{medial\_capture\_check.py}). Remaining errors are ours.}}
\author{}
\date{July 2026}

\begin{document}
\maketitle

\begin{abstract}
Let $\X$ be a closed set in a complete Riemannian manifold $M$. Its
\emph{medial axis} (the cut locus of $\X$) records the points of
ambiguity of the nearest-point problem, and the associated
\emph{reconstruction} is the union of the maximal balls that touch
$\X$ without crossing it. Bia\l{}o\.zyt's Theorem~6 characterises, in
$\R^{n}$, exactly which closed sets are recovered in full by this
construction. We generalise that theorem to \emph{every} complete
Riemannian manifold and \emph{every} nonempty closed proper subset,
with no curvature bound, genericity, analyticity, spread, or
cardinality hypothesis, and we do so from Bia\l{}o\.zyt's theorem
alone.

Two mechanisms drive the theory. The first is a \emph{capture
theorem}. Along any solution of the differential inclusion
$\Phi'\in-\conv U(\Phi)$, where $U$ is the (upper semicontinuous,
generally non-closed-graph-selectable) field of unit directions of
minimising geodesics to $\X$, the distance $\dX$ grows at unit rate
for exactly as long as the trajectory avoids the \emph{two-geodesic
medial axis} $\Mg$. Consequently every point of $M\setminus\X$ is
either \emph{captured} --- it sits at depth $\ge\dX(p)$ inside a
medial ball --- or its trajectory \emph{escapes} to infinity, which
is impossible when $M$ is compact. Hence $\Rg=M\setminus\X$ for every
compact $M$ and every closed $\X\neq\varnothing$, even a single point.
The delicate point --- $\Mg$ is not closed, and the reconstruction
margin is not monotone through multi-geodesic points --- dissolves
once one stops at the \emph{first approach} to $\Mg$ rather than
flowing through it; focal and conjugate degenerations become ordinary
points of the inclusion.

The second mechanism is a horofunction formulation. Calling $\psi$ a
\emph{medial limit function} if $\psi=\lim[d(\cdot,a_{k})-\dX(a_{k})]$
locally uniformly with $a_{k}\in\Mg$, the swallowing of remote balls
becomes a three-line lemma, $\{\psi<0\}\subseteq\Rg$, and
\[
  \Rg=\bigcup\bigl\{\{\psi<0\}:\ \psi\ \text{a medial limit
  function}\bigr\}
\]
holds on every complete $M$. Every point of $M\setminus\X$ is
captured, drift-swallowed, or escaping, and $\Rg$ is exactly the first
two classes. On a Hadamard manifold the ideal medial limit functions
are Busemann functions up to a level, which turns the residual offset
of the older theory into a definition. Bia\l{}o\.zyt's theorem
reappears as the flat specialisation, its supporting-hyperplane
certificate identifying precisely which half-space limits occur; the
compact, spherical, hyperbolic, and Hadamard theorems are all further
specialisations. Throughout we illustrate the definitions and
theorems with explicit examples in $\R^{2}$, on the round sphere, in
the hyperbolic plane, and on surfaces of revolution, and we discuss at
length what the two layers of the theory --- the unconditional
dynamical characterisation and the geometry-specific faithful
certificate --- do and do not deliver.
\end{abstract}

\tableofcontents

\section{Introduction}\label{sec:intro}

\subsection{The reconstruction problem}

Fix a nonempty closed set $\X$ in a metric space and, for each point
$p$ outside $\X$, ask for the nearest points of $\X$. Where the answer
is unique the geometry of $\X$ is locally transparent; the interesting
locus is the set of $p$ with an \emph{ambiguous} nearest-point
answer. In the plane this is the classical \emph{medial axis} or
\emph{skeleton}, a one-dimensional graph that threads the region
between the features of $\X$ and has been used, since Blum, as a
shape descriptor in vision, in mesh generation, and in the topology of
domains \cite{Lie}. Two features of the axis make it a genuine
\emph{descriptor} rather than a mere by-product: each axis point $a$
carries a scalar, the distance $\dX(a)$ to $\X$, and the open ball
$B(a,\dX(a))$ of that radius is a \emph{free ball} --- a maximal ball
that meets $\X$ only in its boundary. The union of the free balls over
the axis,
\[
  \Rx=\bigcup_{a\in\Mx}B\bigl(a,\dX(a)\bigr),
\]
is the \emph{reconstruction} of the complement of $\X$ from the axis.

One wants to know when the reconstruction is faithful, i.e.\ when
$\Rx=M\setminus\X$. The obstruction is easy to see and hard to
characterise. If $\X$ is convex in $\R^{n}$ the axis is empty and
$\Rx=\varnothing$: convex sets cast no medial shadow, so their
complements are not reconstructed at all. If $\X$ is a pair of points,
the axis is the bisecting hyperplane and its free balls fill the two
open half-spaces the pair spans --- but they miss the segment between
the points, which lies in $\conv\X$ and is reconstructed only in the
limit, by balls centred far out along the bisector. The interplay
between the local balls, which fill the ``inside'' of the hull, and
the escaping balls, which fill defective facets of the hull, is the
whole content of the Euclidean theorem.

\subsection{Bia\l{}o\.zyt's theorem, the starting point}

Bia\l{}o\.zyt \cite{Bia} settled the Euclidean case completely. Write
$\conv\X$ for the convex hull and $\overline{\conv}\,\X$ for its
closure. A supporting hyperplane $L$ of $\overline{\conv}\,\X$ is
\emph{defective} if $\X\cap L\neq\overline{\conv}\,\X\cap L$, i.e.\ if
$\X$ fails to fill the contact face that the hull presents on $L$; for
each supporting $L=\{\sca{v}{\cdot}=c\}$ (with $\X\subseteq
\{\sca{v}{\cdot}\le c\}$) let $L^{+}=\{\sca{v}{\cdot}>c\}$ be the open
outer half-space.

\begin{theorem}[Bia\l{}o\.zyt {\cite[Thm.~6]{Bia}}]\label{thm:bia}
Let $\X\subseteq\R^{n}$ be closed and nonempty. Then $\X$ is
reconstructible, i.e.\ $\Rx=\R^{n}\setminus\X$, if and only if
\[
  \R^{n}\setminus\X\ \subseteq\ \overline{\conv}\,\X\ \cup\!\!
  \bigcup_{L\ \text{defective}}\!\!L^{+}.
\]
\end{theorem}

Two features of this statement organise everything that follows.
First, the hull $\overline{\conv}\,\X$ is reconstructed by
\emph{ordinary} (finite) medial balls: every point of $\conv\X$ lies
in a free ball centred on the axis. Second, the outer half-spaces
$L^{+}$ of defective supporting hyperplanes are reconstructed only in
the \emph{limit}, by free balls whose centres escape to infinity along
a direction normal to $L$; a supporting hyperplane contributes such a
limit precisely when its contact set is defective --- equivalently,
by Motzkin's theorem, when a point of $L$ just outside the contact
face is non-Chebyshev. This is the dichotomy \emph{local balls vs.\
escaping balls} in its cleanest form, and the goal of this paper is to
find its counterpart on a curved manifold, where ``half-space'' must
become ``horoball'', ``convex hull'' must become something
curvature-dependent, and the very notion of ``two nearest points''
must be re-examined because two nearest points can be reached by
several competing geodesics.

\subsection{From $\R^{n}$ to a manifold: what breaks}

Let $M$ be a complete Riemannian manifold. Three Euclidean
conveniences fail at once.

\emph{(i) Geodesics are not unique.} Two points can be joined by
several minimising geodesics, and the nearest-point map to $\X$ can be
single-valued (a unique foot $y\in\X$) yet ambiguous as a
\emph{geodesic} problem (several minimisers reaching that same $y$).
The sphere makes this vivid: from any point other than a pole, the
nearest point on a small polar cap is unique, but the antipode of the
cap's centre is reached by a whole circle of minimising geodesics.
This forces a distinction, invisible in $\R^{n}$, between the
\emph{two-point} axis $\Mx$ (two nearest \emph{points}) and the
\emph{two-geodesic} axis $\Mg$ (two nearest \emph{geodesics});
$\Mx\subseteq\Mg$, and it is $\Mg$ --- the cut locus of $\X$ --- that
carries the reconstruction theory.

\emph{(ii) Convexity is not linear.} There is no ``$\conv\X$''; its
role is played by different objects in different curvatures --- the
ordinary hull only in the flat case, a horoconvex hull in
nonpositive curvature, and \emph{nothing at all} on a compact manifold,
where every complement is reconstructed and the hull condition is
vacuous.

\emph{(iii) The distance function turns.} In $\R^{n}$ the ascending
gradient flow of $\dX$ runs along straight rays until it meets the
bisector; on a manifold the minimising geodesic to $\X$ bends, its
foot slides, and the natural ``margin'' one would monotonically
accumulate along the flow genuinely oscillates as the active contact
set changes. This is the analytic heart of the difficulty and the
reason a naive gradient-flow argument stalls exactly at the cut
points, where it is asked to continue past a singularity of $\dX$.

The resolution, developed below, is to change the question the flow is
asked to answer. Rather than flowing \emph{through} the axis and
tracking a monotone quantity, we flow \emph{up to} the axis and read
off a margin that is exact on the complement of the axis and merely
needs to survive to the first approach. This severs the argument from
every regularity subtlety at cut, focal, and conjugate points.

\subsection{Results}

We prove three things and then specialise them.

\emph{Capture (\S\ref{sec:capture}).} Theorem~\ref{thm:capture} shows
that along any trajectory of the descent inclusion, and up to the
first time the trajectory meets $\Mg$, the distance to $\X$ grows at
rate exactly one while the trajectory itself is $1$-Lipschitz, so the
starting point sits at depth $\ge\dX(p)$ inside every ball centred on
the moving trajectory. The trajectory therefore either reaches $\Mg$
in finite time (\emph{capture}: $p\in\Rg$) or escapes to infinity
(\emph{escape}, impossible on a compact manifold). Hence
Corollary~\ref{cor:compact}: $\Rg=M\setminus\X$ on \emph{every}
compact manifold, for every closed $\X\neq\varnothing$ --- including
a single point, which the two-point axis cannot detect at all.

\emph{Drift and the characterisation (\S\ref{sec:drift}).} On a
non-compact manifold, capture is not the whole story: a point may be
reconstructed by remote balls whose centres run to infinity while its
own trajectory escapes. Definition~\ref{def:mlf} packages the limits
of these balls as \emph{medial limit functions}, and
Lemma~\ref{lem:swallow} shows each reconstructs its whole sublevel
set. Theorem~\ref{thm:char} is the resulting characterisation, valid
on every complete manifold, and Theorem~\ref{thm:trichotomy} is the
trichotomy captured / drift-swallowed / escaping.

\emph{Specialisations (\S\ref{sec:spec}).} We recover, as corollaries
and worked examples, the flat theorem (where the abstract
characterisation \emph{is} Bia\l{}o\.zyt's, its half-space limits
identified by the defective-hyperplane certificate), the total
reconstruction of every compact manifold and every spherical space
form, the hyperbolic and Hadamard characterisations (where medial
limit functions are Busemann functions at a level), and the
non-compact positive-curvature examples where reconstruction is
genuinely partial.

Section~\ref{sec:open} isolates what remains open --- entirely on the
\emph{faithful} side, i.e.\ in replacing ``$\psi$ is a medial limit
function'' by a certificate readable from $\X$ and the geometry of
$M$, which is known in the flat and Hadamard cases and open in mixed
curvature. We stress throughout the paper the distinction between the
\emph{unconditional} layer --- the dynamical characterisation, which
this paper proves in full --- and the \emph{faithful} layer --- the
curvature-specific certificate, of which Bia\l{}o\.zyt's theorem is
the flat instance.

\section{Preliminaries}\label{sec:prelim}

Throughout, $M$ is a complete Riemannian manifold; by the
Hopf--Rinow theorem it is a proper geodesic metric space (closed balls
are compact, and any two points are joined by a minimising geodesic).
$\X\subseteq M$ is closed with $\varnothing\neq\X\neq M$, and
$\dX=\dist(\cdot,\X)$ is the distance to $\X$, a $1$-Lipschitz
function, positive exactly off $\X$. All geodesics are parametrised by
arc length.

\subsection{The two axes and the two reconstructions}

For $x\in M\setminus\X$ a \emph{minimising geodesic to $\X$} is a
unit-speed geodesic $\gamma:[0,\dX(x)]\to M$ with $\gamma(0)=x$ and
$\gamma(\dX(x))\in\X$; its foot is $\gamma(\dX(x))$. Set
\[
  U(x)=\bigl\{\gamma'(0):\ \gamma\ \text{a minimising geodesic to
  }\X,\ \gamma(0)=x\bigr\}\subseteq T^{1}_{x}M,
\]
the compact set of unit initial directions of such geodesics (a
geodesic is determined by its initial vector, so distinct minimisers
have distinct directions). Define the two medial axes
\[
  \Mx=\{a:\ a\ \text{has}\ \ge2\ \text{nearest \emph{points} in}\
  \X\},\qquad
  \Mg=\{x\in M\setminus\X:\ |U(x)|\ge2\},
\]
the \emph{two-point} and \emph{two-geodesic} axes. Because two
distinct minimising geodesics may share a foot, $\Mx\subseteq\Mg$;
$\Mg$ is exactly the cut locus of $\X$. Their reconstructions are
\[
  \Rx=\bigcup_{a\in\Mx}B\bigl(a,\dX(a)\bigr),\qquad
  \Rg=\bigcup_{a\in\Mg}B\bigl(a,\dX(a)\bigr),\qquad \Rx\subseteq\Rg.
\]
When we use $\Mg$ as a target to be reached we write $\Sigma:=\Mg$;
$r(a)=\dX(a)$ is the \emph{reach} of a centre $a$. In $\R^{n}$
geodesics are unique, so $\Mx=\Mg$ and $\Rx=\Rg$, and the two theories
coincide; the split is a purely curved phenomenon.

\begin{remark}[Why the two-geodesic axis is the right object]
The two-point axis $\Mx$ is the intuitive skeleton, but it is blind to
a whole class of reconstructions. A single point $\X=\{x_{0}\}$ on the
sphere has $\Mx=\varnothing$ --- every point off $x_{0}$ has the
unique nearest point $x_{0}$ --- yet its antipode $\bar{x}_{0}$ is a
genuine medial centre: the entire sphere minus $x_{0}$ is a free ball
about $\bar{x}_{0}$, because $\bar{x}_{0}$ is joined to $x_{0}$ by a
sphere's worth of minimising geodesics, all of length $\pi$. Thus
$\bar{x}_{0}\in\Mg\setminus\Mx$ and $\Rg=\Sn\setminus\{x_{0}\}$ while
$\Rx=\varnothing$. The reconstruction theory sees this only through
$\Mg$. This example, trivial as it is, already shows that the compact
theorem below \emph{cannot} hold for $\Mx$ (a single point is never
reconstructed by two-point balls) and \emph{must} be stated for
$\Mg$.
\end{remark}

\subsection{The direction field}

\begin{lemma}[Direction field: compactness and semicontinuity]
\label{lem:usc}
For every $x\in M\setminus\X$ the set $U(x)$ is nonempty and compact,
and the multimap $x\mapsto U(x)$ is upper semicontinuous: if
$x_{k}\to x$ in $M\setminus\X$ and $u_{k}\in U(x_{k})$ with
$u_{k}\to u$ in $TM$, then $u\in U(x)$. In particular $U$ is
single-valued and continuous at every point of
$(M\setminus\X)\setminus\Sigma$.
\end{lemma}

\begin{proof}
Nonemptiness and compactness: minimisers exist by properness, and a
limit of initial vectors of minimising geodesics from a fixed $x$ is
the initial vector of a limit geodesic, minimising by continuity of
$\dX$ and of length. Upper semicontinuity: let $\gamma_{k}$ be the
minimising geodesic with $\gamma_{k}'(0)=u_{k}$. By smooth dependence
of geodesics on initial conditions, $\gamma_{k}\to\gamma$ locally
uniformly, where $\gamma'(0)=u$; the lengths $\dX(x_{k})\to\dX(x)$ (as
$\dX$ is continuous), and the feet
$\gamma_{k}(\dX(x_{k}))\in\X$ converge to $\gamma(\dX(x))$, which lies
in $\X$ since $\X$ is closed. So $\gamma$ is a minimising geodesic from
$x$ to $\X$, i.e.\ $u\in U(x)$. At a point with $U(x)=\{u\}$, upper
semicontinuity of a compact-valued map into a singleton forces
continuity of every selection \emph{at} $x$ --- though $U$ may be
multivalued arbitrarily close by.
\end{proof}

\begin{lemma}[First variation of $\dX$; {\cite[\S3]{MM}}]
\label{lem:firstvar}
For $x\in M\setminus\X$ and $w\in T_{x}M$,
\[
  \frac{d}{ds}\Big|_{s=0^{+}}\dX\bigl(\exp_{x}(sw)\bigr)
  =\min_{u\in U(x)}\sca{w}{-u}.
\]
\end{lemma}

\begin{proof}[Proof sketch]
``$\le$'': fix $u\in U(x)$ with foot $y$; comparing $\dX$ with the
distance to $y$ along the varied geodesics gives
$\dX(\exp_{x}(sw))\le d(\exp_{x}(sw),y)\le\dX(x)+s\sca{w}{-u}+O(s^{2})$,
by the first variation of arc length. ``$\ge$'': for $s_{k}\downarrow0$
pick $u_{k}\in U(\exp_{x}(s_{k}w))$; a subsequence converges to some
$u\in U(x)$ (Lemma~\ref{lem:usc}), and the comparison along the
geodesics with directions $u_{k}$ gives the matching bound with this
$u$, hence with the minimum. The semiconcavity of $\dX$ on
$M\setminus\X$, which we shall not need, and full details, are in
Mantegazza--Mennucci \cite{MM}.
\end{proof}

Intuitively: $\dX$ increases fastest, at unit rate, in the direction
directly \emph{away} from a nearest foot; at a two-geodesic point the
rate in a given direction $w$ is dictated by the \emph{least
favourable} competing foot, which is what makes the min appear and
what makes the gradient flow turn.

\subsection{The descent inclusion and capture curves}

On the open set $M\setminus\X$ define the multimap
\[
  G(x)\ :=\ \conv\bigl(-U(x)\bigr)\ \subseteq\ T_{x}M ,
\]
the convex hull of the reversed minimising directions. By
Lemma~\ref{lem:usc}, $G$ is upper semicontinuous with nonempty compact
convex values contained in the closed unit ball. A \emph{capture
curve} from $p\in M\setminus\X$ is a locally Lipschitz curve
$\Phi:[0,T_{\max})\to M\setminus\X$ with $\Phi_{0}=p$ and
\begin{equation}\label{eq:incl}
  \Phi'(t)\in G(\Phi_{t})\qquad\text{for a.e.\ }t ,
\end{equation}
maximally extended in $M\setminus\X$. Existence of a solution, and its
extension for as long as the trajectory stays in a compact subset of
the open domain $M\setminus\X$, are the standard theory of
differential inclusions with upper semicontinuous, convex, compact
right-hand side (Aubin--Cellina \cite[\S2.1]{AC}), applied in charts
and concatenated; since $G(x)$ lies in the unit ball, $|\Phi'|\le1$
a.e.\ and trajectories are $1$-Lipschitz. No uniqueness, canonicity,
or semiconcave-gradient selection is required: \emph{any} solution
serves. Off $\Sigma$ the field $-U$ is single-valued
(Lemma~\ref{lem:usc}) and \eqref{eq:incl} is the classical ODE
$\Phi'=-u(\Phi)$ --- the curve simply extends the current minimising
geodesic backwards from $x$, its foot sliding continuously; at a point
of $\Sigma$ the inclusion may branch, and we shall never analyse it
there.

\begin{remark}[Relation to Lieutier's flow]\label{rem:lieutier}
For semiconcave $\dX$ one can canonically select the minimal-norm
element of $G$, giving Lieutier's generalised-gradient flow \cite{Lie}
(in $\R^{n}$; \cite{ACa,MM} for the semiconcavity). Our argument uses
the \emph{opposite} regime: it analyses \eqref{eq:incl} only
\emph{off} $\Sigma$, where every selection agrees with the naive one,
and it \emph{terminates} at $\Sigma$, where Lieutier's flow is
designed to keep going. This is exactly why no analogue of the swallow
property \emph{past} cut points is ever needed --- the classical
sticking point in extending the Euclidean argument to manifolds.
\end{remark}

\subsection{Free-ball functions and medial limit functions}

For $a\in M\setminus\X$ set
\[
  \psi_{a}:=d(\cdot,a)-\dX(a)\ :\ M\to\R ,
\]
a $1$-Lipschitz function, negative exactly on the open free ball
$B(a,\dX(a))$ and $\ge0$ on $\X$. Crucially it is normalised by the
\emph{reach}: $\min\psi_{a}=-\dX(a)$, attained at $a$, with no
reference basepoint. This intrinsic normalisation --- ``measure the
ball by how deep it goes, not by how far its centre is'' --- is what
makes the limits below basepoint-free, and it is the technical device
that dissolves the offset bookkeeping of the older Hadamard theory.

\begin{definition}[Medial limit functions]\label{def:mlf}
A function $\psi:M\to\R$ is a \emph{medial limit function} of $\X$ if
there exist $a_{k}\in\Mg$ with $\psi_{a_{k}}\to\psi$ locally
uniformly. It is \emph{finite} if $(a_{k})$ has a bounded subsequence
and \emph{ideal} otherwise (then $d(\cdot,a_{k})\to\infty$ locally
uniformly while $\psi$ stays finite, so $\psi$ is of horofunction
type). Each such $\psi$ is $1$-Lipschitz and $\ge0$ on $\X$, so
$\{\psi<0\}$ is an open set disjoint from $\X$: the \emph{free
horoball} of $\psi$. For a finite $\psi$ along $a_{k}\to a$: if
$a\in\Mg$ then $\psi=\psi_{a}$ is an honest free-ball function; in
general ($\Mg$ is not closed) $\psi=\psi_{a}$ with $a$ a
focal-degenerate limit of medial centres --- still usable by
Lemma~\ref{lem:swallow}, which never needs the limit centre itself.
\end{definition}

\begin{remark}[Relation to the horofunction compactification]
\label{rem:horo}
Since $M$ is proper, the space $C(M)$ with locally uniform convergence
is a compact home for $1$-Lipschitz functions once they are normalised
(Arzel\`a--Ascoli). The classical horofunction compactification
\cite{Gro,BH} normalises $d(\cdot,a_{k})$ by subtracting its value at
a fixed basepoint $o$; our $\psi_{a_{k}}$ instead subtracts $\dX(a_{k})$,
so the two differ by the scalar $d(o,a_{k})-\dX(a_{k})$ --- the
\emph{offset}. Reach-normalisation absorbs the offset into the
function; see Remark~\ref{rem:offset}. On a Hadamard manifold every
horofunction is a Busemann function (horofunctions of a proper
$\mathrm{CAT}(0)$ space are Busemann, \cite[II.8.13]{BH}), so every
ideal medial limit function is ``$b_{\xi}+c$'' for some
$\xi\in\bdry M$ and level $c\in\R$.
\end{remark}

\section{The capture theorem}\label{sec:capture}

The capture theorem is the engine of the compact case and one of the
two inputs to the general characterisation. Its statement is a clean
dichotomy; its proof is a one-line differential identity valid off the
axis, plus a compactness argument for the escape alternative.

\begin{theorem}[Capture or escape]\label{thm:capture}
Let $M$ be complete, $\X$ closed, $p\in M\setminus\X$, and let $\Phi$
be any capture curve from $p$. Set
\[
  T\ :=\ \sup\bigl\{t\in[0,T_{\max}):\ \Phi_{s}\notin\Sigma\ \
  \forall s\in[0,t]\bigr\}.
\]
Then on $[0,T)$,
\begin{equation}\label{eq:rateone}
  \dX(\Phi_{t})=\dX(p)+t,
  \qquad
  d(p,\Phi_{t})\le t,
  \qquad\text{hence}\qquad
  \dX(\Phi_{t})-d(p,\Phi_{t})\ \ge\ \dX(p),
\end{equation}
and exactly one of the following holds.
\begin{enumerate}[label=\textup{(\roman*)}]
\item\emph{(Capture.)} $T<T_{\max}$. Then for every $\varepsilon>0$
there is $a\in\Mg$ (with $a=\Phi_{T}$ itself if $\Phi_{T}\in\Sigma$)
such that
\[
  \dX(a)-d(p,a)\ \ge\ \dX(p)-\varepsilon,
\]
so $p\in B(a,\dX(a))$: hence $p\in\Rg$, at depth arbitrarily close to
$\dX(p)$, and $T\le\operatorname{diam}(M)-\dX(p)$.
\item\emph{(Escape.)} $T=T_{\max}$. Then $T_{\max}=\infty$, the
trajectory never meets $\Sigma$, $\dX(\Phi_{t})=\dX(p)+t\to\infty$, and
$\Phi_{t}$ leaves every compact subset of $M$. This is possible only
if $M$ is non-compact.
\end{enumerate}
\end{theorem}

\begin{proof}
\emph{The rate-one identity.} Fix $t\in[0,T)$; on $[0,t]$ the
trajectory avoids $\Sigma$, so $U(\Phi_{s})=\{u(\Phi_{s})\}$ is a
singleton for every $s\le t$ and \eqref{eq:incl} reads
$\Phi'(s)=-u(\Phi_{s})$ at a.e.\ $s$. Put $h(s):=\dX(\Phi_{s})$,
a Lipschitz function. Let $s$ be a point where $\Phi'(s)$ exists
(satisfying the inclusion) and $h'(s)$ exists --- almost every
$s\in[0,t]$. Write $v=\Phi'(s)=-u(\Phi_{s})$, a unit vector. In a chart
around $\Phi_{s}$, differentiability of $\Phi$ gives
$d(\Phi_{s+\sigma},\exp_{\Phi_{s}}(\sigma v))=o(\sigma)$, and $\dX$ is
$1$-Lipschitz, so by Lemma~\ref{lem:firstvar} and single-valuedness,
\[
  h(s+\sigma)=\dX\bigl(\exp_{\Phi_{s}}(\sigma v)\bigr)+o(\sigma)
  =h(s)+\sigma\min_{u\in U(\Phi_{s})}\sca{v}{-u}+o(\sigma)
  =h(s)+\sigma\,\sca{-u}{-u}+o(\sigma).
\]
Thus $h'(s)=|u|^{2}=1$ a.e.\ on $[0,t]$, and integrating,
$h(t)=\dX(p)+t$. The bound $d(p,\Phi_{t})\le t$ is the $1$-Lipschitz
property of $\Phi$, and the margin inequality is their difference. No
closedness of $\Sigma$ and no structure at focal or conjugate points
was used: a point of $(M\setminus\X)\setminus\Sigma$ in the closure of
$\Sigma$ is simply a single-valued point of $U$, and the identity
passes straight through it.

\emph{Dichotomy.} Suppose $T<T_{\max}$. If $\Phi_{T}\in\Sigma$, then by
continuity of the margin and \eqref{eq:rateone} on $[0,T)$,
$\dX(\Phi_{T})-d(p,\Phi_{T})\ge\dX(p)>0$, so $a=\Phi_{T}\in\Mg$
contains $p$ in its open medial ball at depth $\ge\dX(p)$. If
$\Phi_{T}\notin\Sigma$, then by the definition of $T$ as a supremum
every interval $(T,T+\delta]$ contains a time $s_{\delta}$ with
$\Phi_{s_{\delta}}\in\Sigma$; letting $\delta\downarrow0$ and using
continuity of the margin, $a=\Phi_{s_{\delta}}\in\Mg$ satisfies
$\dX(a)-d(p,a)\ge\dX(p)-\varepsilon$ for $\delta$ small. Either way
$p\in B(a,\dX(a))$ with $a\in\Mg$, so $p\in\Rg$; and
$\dX\le\operatorname{diam}(M)$ bounds $T$ through \eqref{eq:rateone}.

If instead $T=T_{\max}$, the trajectory never meets $\Sigma$ and
\eqref{eq:rateone} holds on all of $[0,T_{\max})$. Were
$T_{\max}<\infty$, then $\Phi([0,T_{\max}))\subseteq\bar
B(p,T_{\max})\cap\{\dX\ge\dX(p)\}$, a compact subset of the open set
$M\setminus\X$ (properness), and the extension property would
contradict maximality. So $T_{\max}=\infty$ and
$\dX(\Phi_{t})=\dX(p)+t\to\infty$; as $\dX$ is bounded on compact sets,
$\Phi_{t}$ leaves every compact set. On a compact $M$,
$\dX\le\operatorname{diam}(M)$ forbids this, forcing case (i).
\end{proof}

\begin{corollary}[Total reconstruction on compact manifolds]
\label{cor:compact}
Let $M$ be a compact Riemannian manifold and $\X\subseteq M$ closed,
$\varnothing\neq\X\neq M$. Then
\[
  \boxed{\ \Rg=M\setminus\X\ },
\]
and every $p\in M\setminus\X$ lies at depth $\ge\dX(p)-\varepsilon$
(any $\varepsilon>0$) inside a medial ball. No curvature, genericity,
analyticity, spread, or cardinality hypothesis is required; in
particular $|\X|=1$ is allowed, where $\Mx=\varnothing$ but $\Mg$
reconstructs everything.
\end{corollary}

\begin{proof}
$\Rg\subseteq M\setminus\X$ because medial balls are free. Conversely,
if $p\in\Mg$ then $p\in B(p,\dX(p))\subseteq\Rg$; otherwise start a
capture curve at $p$ (local existence, \S\ref{sec:prelim}). Escape is
impossible on compact $M$ (Theorem~\ref{thm:capture}), so $p$ is
captured, $p\in\Rg$.
\end{proof}

\begin{corollary}[Two-point axis under spread]\label{cor:spread}
If $M$ is compact and $\dX<\inj(M)$ on $M\setminus\X$ --- e.g.\ if
$\X$ is $\delta$-dense with $\delta<\inj(M)$ --- then
$\Rx=M\setminus\X$ for the two-point axis as well.
\end{corollary}

\begin{proof}
By Corollary~\ref{cor:compact} every $p$ lies in $B(a,\dX(a))$ for
some $a\in\Mg$. Two distinct minimising geodesics from $a$ with a
common foot $x_{0}$ would place $x_{0}$ in the cut locus of $a$, giving
$\dX(a)=d(a,x_{0})\ge\inj(M)$, a contradiction. So the two geodesics of
$a$ have distinct feet: $a\in\Mx$.
\end{proof}

\begin{remark}[Why stopping \emph{at} $\Sigma$ is the correct move]
\label{rem:notmonotone}
Two tempting strengthenings of Theorem~\ref{thm:capture} are false,
and their failure is what shapes the proof.

\emph{The linear margin is not monotone through multi-geodesic arcs},
already in $\R^{2}$. Take $\X=\{(0,0),(0,2)\}$ and $p=(1,0.9)$. The
Lieutier flow runs to the bisector $\{y=1\}$ and slides along it toward
$+\infty$ in $x$; the margin $\dX(\Phi)-d(p,\Phi)$ \emph{decreases}
strictly, from $\approx1.345$ toward its limit $1$ (the depth of $p$
in the limiting half-plane $\{x>0\}$'s bisector ball). So one cannot
hope to \emph{accumulate} margin by flowing through the axis.

\emph{Even the quadratic margin} $\dX^{2}-d(p,\cdot)^{2}$, which
\emph{is} monotone along unique-foot and along bisector arcs in
$\R^{n}$, fails at turning points where the active contact set
changes. Three points suffice: $\X=\{(0,0),(0,2),(10,5)\}$, with the
flow from $p=(1,0.9)$ turning at $\approx(5.75,1)$, where the
derivative of the quadratic margin is negative.

Both phenomena live \emph{on} $\Sigma$; on its complement the rate-one
identity \eqref{eq:rateone} is exact. Since reaching $\Sigma$ is the
\emph{goal}, not an obstacle, the correct move is to stop there ---
capture is achieved, with the margin intact, \emph{before} the first
multi-geodesic time. The failure of closedness of $\Sigma$ (focal
degenerations: the evolute of an ellipse in $\R^{2}$, or the
conjugate-cusp endpoints of an ellipsoid's cut segment) is equally
harmless: such points are unique-geodesic points, $-U$ is continuous
there, and the trajectory passes through with $\dX$ growing at rate
one, the foot sliding continuously.
\end{remark}

\section{Drift: the swallow, the characterisation, the trichotomy}
\label{sec:drift}

Capture settles the compact case, but on a non-compact $M$ a point can
be reconstructed by medial balls whose centres escape to infinity while
its own capture curve escapes as well. The clean example is
$\X=\{x_{0},x_{1}\}\subseteq\R^{2}$: a point $p$ in an open half-plane
bounded by the line $x_{0}x_{1}$, placed beyond $x_{0}$ at an obtuse
angle, has a capture curve running straight to infinity with the
single foot $x_{0}$ --- yet $p$ is swallowed by the bisector balls
$B(a_{k},r(a_{k}))$ with $a_{k}\to\infty$, whose limit is a half-plane.
The right general home for such limits is the space of medial limit
functions (Definition~\ref{def:mlf}).

\begin{lemma}[Swallow]\label{lem:swallow}
Let $M$ be any complete Riemannian manifold, $\X$ closed, and $\psi$ a
medial limit function of $\X$. Then $\{\psi<0\}\subseteq\Rg$. If the
witnessing centres lie in $\Mx$, then $\{\psi<0\}\subseteq\Rx$.
\end{lemma}

\begin{proof}
Let $\psi=\lim\psi_{a_{k}}$ locally uniformly, $a_{k}\in\Mg$, and let
$q\in\{\psi<0\}$. Then $\psi_{a_{k}}(q)\to\psi(q)<0$, so for large $k$,
$d(q,a_{k})-\dX(a_{k})=\psi_{a_{k}}(q)<0$, i.e.\ $q\in B(a_{k},\dX(a_{k}))$
with $a_{k}\in\Mg$: $q\in\Rg$.
\end{proof}

There is no anchoring of feet, no curvature bound, no pinching, no
offset hypothesis: the entire content of the swallowing step ---
which in curved settings historically required delicate control of
where the escaping centres' feet accumulate --- is absorbed into the
normalisation of $\psi_{a}$ by the reach.

\begin{lemma}[Compactness of reconstructing sequences]\label{lem:cpt}
Let $p\in\Rg$, and let $a_{k}\in\Mg$ satisfy
$\limsup_{k}[\dX(a_{k})-d(p,a_{k})]>0$ (for instance, centres whose
balls contain $p$ with margin bounded away from $0$; a single repeated
ball qualifies). Then a subsequence of $(\psi_{a_{k}})$ converges
locally uniformly to a medial limit function $\psi$ with $\psi(p)<0$.
\end{lemma}

\begin{proof}
The $\psi_{a_{k}}$ are $1$-Lipschitz and
$\psi_{a_{k}}(p)=d(p,a_{k})-\dX(a_{k})$ is bounded: below by $-\dX(p)$
(since $\dX(a_{k})\le\dX(p)+d(p,a_{k})$) and above by $0$ along the
qualifying subsequence. By Arzel\`a--Ascoli on the proper space $M$, a
further subsequence converges locally uniformly to a medial limit
function $\psi$, and $\psi(p)\le-\limsup[\dX(a_{k})-d(p,a_{k})]<0$.
\end{proof}

\begin{theorem}[Characterisation on an arbitrary complete manifold]
\label{thm:char}
For every complete Riemannian manifold $M$ and every closed
$\X\subseteq M$, $\varnothing\neq\X\neq M$,
\[
  \boxed{\ \Rg\ =\ \bigcup\bigl\{\{\psi<0\}:\ \psi\ \text{a medial
  limit function of}\ \X\bigr\}\ }.
\]
The same identity holds for $\Rx$ and $\Mx$ in place of $\Rg$ and
$\Mg$.
\end{theorem}

\begin{proof}
``$\supseteq$'' is Lemma~\ref{lem:swallow}. ``$\subseteq$'': if
$p\in\Rg$ then $p\in B(a,\dX(a))$ for some $a\in\Mg$, and the constant
sequence $a_{k}\equiv a$ makes $\psi_{a}$ a (finite) medial limit
function with $\psi_{a}(p)<0$.
\end{proof}

The identity is deliberately economical, and one should read its force
correctly. Its finite members alone give a tautology --- every ball
is its own limit --- so the set equality is not the point. The point
is the pair of facts it packages. \emph{First}, the \emph{ideal}
members reconstruct their entire open sublevel set, bare, with no
anchoring or pinching hypothesis: this is the content of
Lemma~\ref{lem:swallow} applied to horofunction-type limits.
\emph{Second}, combined with the capture theorem it yields a clean
trichotomy.

\begin{theorem}[Trichotomy on a complete manifold]
\label{thm:trichotomy}
Let $M$ be complete, $\X$ closed, $p\in M\setminus\X$. Exactly one of
the following holds.
\begin{enumerate}[label=\textup{(\arabic*)}]
\item\emph{(Captured.)} Some capture curve from $p$ approaches $\Mg$
in finite time. Then $p\in\Rg$, at depth arbitrarily close to
$\dX(p)$.
\item\emph{(Swallowed, uncaptured.)} No capture curve from $p$
approaches $\Mg$ in finite time, but $\psi(p)<0$ for some medial limit
function $\psi$. Then $p\in\Rg$: reconstructed, but only by centres
$p$'s own ascent cannot find --- a remote ball, or a drifting
sequence.
\item\emph{(Escaping.)} Neither. Then $p\notin\Rg$, and every capture
curve from $p$ escapes to infinity with $\dX(\Phi_{t})=\dX(p)+t$.
\end{enumerate}
In particular $\Rg=\{p:\ \textup{(1) or (2)}\}$, and on compact $M$
class \textup{(1)} is all of $M\setminus\X$.
\end{theorem}

\begin{proof}
The classes are exclusive and exhaustive by construction. (1)
$\Rightarrow p\in\Rg$ is Theorem~\ref{thm:capture}(i); (2)
$\Rightarrow p\in\Rg$ is Lemma~\ref{lem:swallow}. Conversely, if
$p\in\Rg$ then $\psi_{a}(p)<0$ for the finite medial limit function of
some containing ball, so $p$ lies in class (2) whenever not in class
(1) --- never in (3). A point outside class (1) has every capture
curve uncaptured, hence escaping, by Theorem~\ref{thm:capture}.
\end{proof}

Bia\l{}o\.zyt's theorem is not a description of the reconstructed set
but a \emph{reconstructibility criterion}: it says exactly when a
closed set is recovered in full. The trichotomy yields the direct
analogue on any complete manifold, because $\Rg$ is precisely classes
(1)--(2) while class (3) is its complement.

\begin{corollary}[Reconstructibility criterion]
\label{cor:reconstructible}
Let $M$ be complete, $\X$ closed, $\varnothing\neq\X\neq M$. The
following are equivalent.
\begin{enumerate}[label=\textup{(\alph*)}]
\item $\X$ is reconstructible: $\Rg=M\setminus\X$.
\item $M\setminus\X\ \subseteq\ \bigcup\bigl\{\{\psi<0\}:\ \psi\
\text{a medial limit function of}\ \X\bigr\}$.
\item Every $p\in M\setminus\X$ is captured or drift-swallowed;
equivalently, the escaping class \textup{(3)} of
Theorem~\ref{thm:trichotomy} is empty.
\item For every $p\in M\setminus\X$ whose capture curves escape to
infinity, some medial limit function $\psi$ satisfies $\psi(p)<0$.
\end{enumerate}
\end{corollary}

\begin{proof}
$(a)\Leftrightarrow(b)$: by Theorem~\ref{thm:char} the right-hand side
of (b) equals $\Rg$, and $\Rg\subseteq M\setminus\X$ always, so (b) is
the reverse inclusion, i.e.\ (a). $(a)\Leftrightarrow(c)$: by
Theorem~\ref{thm:trichotomy}, $\Rg$ is exactly classes (1)--(2) and is
disjoint from class (3); since the three classes partition
$M\setminus\X$, $\Rg=M\setminus\X$ iff class (3) is empty.
$(c)\Leftrightarrow(d)$: a point outside class (1) escapes
(Theorem~\ref{thm:capture}), so it lies in class (2) --- swallowed ---
iff some $\psi$ has $\psi(p)<0$; captured points lie in $\Rg$
regardless, so emptiness of class (3) is exactly (d).
\end{proof}

Two remarks fix its meaning. First, \emph{on a compact manifold the
criterion is vacuously satisfied}: escape is impossible
(Theorem~\ref{thm:capture}), class (3) is empty, and $\X$ is
reconstructible for free (Corollary~\ref{cor:compact}). The criterion
therefore has content only on non-compact $M$ --- exactly where
Bia\l{}o\.zyt's does, $\R^{n}$ being non-compact. Second, condition
(b) is the \emph{unconditional} form: it refers to the medial limit
functions, hence implicitly to the axis $\Mg$. Bia\l{}o\.zyt's theorem
is the \emph{faithful} sharpening of (b) in the flat case, replacing
``$M\setminus\X\subseteq\bigcup\{\psi<0\}$'' by the axis-free hull
condition
\[
  \R^{n}\setminus\X\ \subseteq\ \overline{\conv}\,\X\ \cup\!\!
  \bigcup_{L\ \text{defective}}\!\!L^{+}
\]
of Theorem~\ref{thm:bia} (Corollary~\ref{cor:flat}); on a Hadamard
manifold the corresponding faithful sharpening replaces the ideal
sublevel sets by defective horoballs $\{b_{\xi}<-c(\xi)\}$
(\S\ref{sec:hadamard}, Remark~\ref{rem:offset}). Thus
Corollary~\ref{cor:reconstructible} is the general analogue of
Bia\l{}o\.zyt's main theorem, and the passage from (b) to a
checkable-from-$\X$ certificate is precisely the faithful-layer
problem, solved in flat and Hadamard geometry and open in mixed
curvature.

\begin{remark}[The witnesses of class (2)]\label{rem:class2}
Both kinds of witness genuinely occur, already for
$\X=\{x_{0},x_{1}\}\subseteq\R^{2}$ and $p$ beyond $x_{0}$ in an open
half-plane at an obtuse angle to the bisector. Every capture curve from
$p$ runs straight to infinity with the single foot $x_{0}$, yet $p$
lies in a single \emph{finite} bisector ball (and in the ideal
half-plane limit of the bisector balls). Class (2) is thus the precise
sense in which ``reconstruction is non-local for $p$'': the
certificate cannot be found by ascending from $p$. This is exactly the
phenomenon that separates the non-compact theory from the compact one,
and the reason the faithful certificates below must quantify over
supporting half-spaces or horoballs rather than over flow lines.
\end{remark}

\begin{remark}[The Hadamard offset, dissolved]\label{rem:offset}
Let $M$ be Hadamard (complete, simply connected, $\sect\le0$). By
Remark~\ref{rem:horo} the ideal medial limit functions are exactly
$b_{\xi}+c$ with $\xi\in\bdry M$ and $c\in\R$, normalised so that
$b_{\xi}$ vanishes on the horosphere of the maximal free horoball
$\Hb{}$ at $\xi$; necessarily $c\ge0$ when $\X$ meets every
neighbourhood of that horosphere. Call $\Hb{}$ \emph{defective at level
$c$} if $b_{\xi}+c$ is a medial limit function; then
Lemma~\ref{lem:swallow} reconstructs $\{b_{\xi}<-c\}$, the horoball
shrunk by $c$. The historical ``radial offset''
$\mathrm{off}_{k}=d(o,a_{k})+b_{\xi}(a_{k})$ is precisely the
discrepancy $\psi_{a_{k}}-b_{\xi}$ along the drift, so a drift with
$\mathrm{off}_{k}\to c$ is a level-$c$ certificate. What used to be a
gap --- ``only the deep part $\{b_{\xi}<-c\}$ is reconstructed for a
non-vanishing asymptotic offset'' --- is now simply the content of the
level: the drift \emph{is} a level-$c$ certificate and reconstructs
exactly $\{b_{\xi}<-c\}$. The full-horoball statement is the assertion
that a level-$0$ certificate exists, and the exact level
\[
  c(\xi)\ :=\ \inf\{c\ge0:\ b_{\xi}+c\ \text{is a medial limit
  function}\}
\]
is a question about $\X$ and the horosphere geometry, not about the
medial axis; see \S\ref{sec:open}.
\end{remark}

\section{Specialisations and worked examples}\label{sec:spec}

We now recover the classical theory as specialisations of
Theorems~\ref{thm:char}--\ref{thm:trichotomy}, and illustrate every
key definition with an explicit example. Throughout, the
\emph{unconditional} content --- the abstract characterisation --- is
supplied by the theorems above; the \emph{faithful} content --- a
certificate readable from $\X$ and $M$ --- is supplied by
Bia\l{}o\.zyt's theorem in the flat case and discussed as the open
residue elsewhere.

\subsection{Euclidean space: recovering Bia\l{}o\.zyt's theorem}
\label{sec:euclid}

In $\R^{n}$ geodesics are unique, so $\Mg=\Mx=:M_{\X}$ and
$\Rg=\Rx=:\Rx$. We identify the medial limit functions. A finite
medial limit function is a free-ball function $y\mapsto|y-a|-r$,
$a\in\overline{M_{\X}}$. For an \emph{ideal} one, let $a_{k}\to\infty$
along a unit direction $u$ (i.e.\ $a_{k}/|a_{k}|\to u$) with reaches
$r_{k}=\dX(a_{k})$. Then, for each fixed $y$,
\[
  \psi_{a_{k}}(y)=|y-a_{k}|-r_{k}
  =\bigl(|y-a_{k}|-|a_{k}|\bigr)+\bigl(|a_{k}|-r_{k}\bigr)
  \longrightarrow\ -\sca{y}{u}+\beta,
\]
where the first bracket tends to $-\sca{y}{u}$ (a standard Euclidean
computation) and convergence of $\psi_{a_{k}}$ forces the offset
$|a_{k}|-r_{k}$ to a finite limit $\beta$. So every ideal medial limit
function is \emph{affine},
\[
  \psi(y)=\beta-\sca{y}{u},\qquad\{\psi<0\}=\{\sca{y}{u}>\beta\},
\]
an open half-space with outward normal $u$. Conversely a supporting
half-space limit of medial balls is exactly such a $\psi$.
Theorem~\ref{thm:char} therefore reads, in $\R^{n}$:
\[
  \Rx\ =\ \Bigl(\bigcup\text{finite medial balls}\Bigr)\ \cup\
  \Bigl(\bigcup\bigl\{\text{open half-spaces that are ideal medial
  limits}\bigr\}\Bigr).
\]
This is Bia\l{}o\.zyt's decomposition. The finite balls fill
$\conv\X\setminus\X$ (every point of the convex hull lies in a free
ball); the half-spaces are the outer half-spaces $L^{+}$ of
supporting hyperplanes $L$ that arise as ideal limits. Bia\l{}o\.zyt's
Theorem~\ref{thm:bia} identifies \emph{which} supporting hyperplanes
occur: precisely the \emph{defective} ones, $\X\cap
L\neq\overline{\conv}\,\X\cap L$ (equivalently, by Motzkin's theorem,
those carrying a non-Chebyshev point just outside the contact face).
Thus:

\begin{corollary}[Flat specialisation]\label{cor:flat}
In $\R^{n}$, Theorem~\ref{thm:char} together with
Bia\l{}o\.zyt's Theorem~\ref{thm:bia} gives
\[
  \Rx\ =\ \bigl(\overline{\conv}\,\X\setminus\X\bigr)\ \cup\!\!
  \bigcup_{L\ \text{defective}}\!\!L^{+},
\]
and $\X$ is reconstructible iff
$\R^{n}\setminus\X\subseteq\overline{\conv}\,\X\cup\bigcup_{L}L^{+}$.
The drift over a defective $L$ runs at level $\beta=0$ (no offset): the
escaping bisector centres accumulate their feet on the defective
contact face, and their half-space limit is $L^{+}$ itself.
\end{corollary}

Here the split of labour is exact and instructive: the abstract
characterisation (unconditional, self-contained) delivers ``$\Rx$ is
balls $\cup$ half-space limits''; Bia\l{}o\.zyt's theorem (the faithful
certificate) delivers ``the half-space limits are exactly the
defective supporting hyperplanes.'' This is the template the whole
paper follows on other geometries.

\begin{example}[Two points]\label{ex:twopoints}
$\X=\{x_{0},x_{1}\}\subseteq\R^{2}$, so $\overline{\conv}\,\X$ is the
segment $[x_{0},x_{1}]$ and the axis $M_{\X}$ is its perpendicular
bisector. There is exactly one defective supporting line: the line
$\ell$ \emph{through} $x_{0}x_{1}$ itself, on which
$\overline{\conv}\,\X\cap\ell=[x_{0},x_{1}]$ while $\X\cap\ell=
\{x_{0},x_{1}\}$ --- two points, not the whole segment. It supports
from both sides, so it is defective in both orientations, with the two
open half-planes bounded by $\ell$ as outer half-spaces $L^{+}$; these
are the ideal (half-plane) medial limits of the escaping bisector
balls. By Corollary~\ref{cor:flat},
\[
  \Rx=\bigl([x_{0},x_{1}]\setminus\{x_{0},x_{1}\}\bigr)\ \cup\
  \{\text{two open half-planes}\}\ =\
  \R^{2}\setminus\bigl(\ell\setminus(x_{0},x_{1})\bigr),
\]
i.e.\ \emph{everything except the two outer rays} of $\ell$ beyond the
endpoints (together with $\X$ itself). Note what is and is not
reconstructed: the \emph{open} segment $(x_{0},x_{1})$ \emph{is}
reconstructed --- by finite bisector balls, which contain every
interior segment point --- while the two collinear rays $\{x_{0}+
t(x_{0}-x_{1}):t>0\}$ and its mirror are on \emph{no} free ball and are
missed by the open half-planes. Removing the reconstruction of these
rays is exactly the failure of reconstructibility for a two-point set.
This is the primal drift example of Remark~\ref{rem:class2}: a point in
an open half-plane at an obtuse angle to the bisector is reconstructed
only by the ideal half-plane limit (class (2)), not by its own escaping
capture curve.
\end{example}

\begin{example}[A triangle: the outer half-spaces close up]
\label{ex:arc}
Let $\X$ be the three vertices of an equilateral triangle, so
$\overline{\conv}\,\X$ is the solid triangle. A supporting line through
a single vertex has $\X\cap L=\overline{\conv}\,\X\cap L=\{v\}$ ---
not defective. A supporting line \emph{along an edge} has
$\overline{\conv}\,\X\cap L$ the whole edge while $\X\cap L$ is its two
endpoints --- defective, with a genuinely nonempty outer half-plane
$L^{+}$ reconstructed by the ideal drift of the escaping bisector
centres of those two vertices (which are legitimate medial centres far
out beyond the edge, the third vertex being strictly farther). The
three edge half-planes \emph{overlap to cover every exterior point} ---
a point beyond a vertex lies in the outer half-plane of an incident
edge --- and the interior is filled by finite balls, so
$\Rx=\R^{2}\setminus\X$: reconstructible. Contrast
Example~\ref{ex:twopoints}, the degenerate ``triangle'' with one edge:
there the single edge's two outer half-planes leave the collinear rays
uncovered, and reconstructibility fails. The moral is that
reconstructibility is about whether the defective outer half-spaces
\emph{tile} the exterior, which a genuine polygon does and a segment
does not.
\end{example}

\subsection{Compact manifolds, the sphere, and space forms}
\label{sec:compact}

Here the faithful layer is \emph{vacuous}: there are no ideal medial
limit functions (the manifold is bounded), no hull condition, and
Corollary~\ref{cor:compact} gives total reconstruction outright.

\begin{corollary}[Round sphere and spherical space forms]
\label{cor:sphere}
For every closed $\varnothing\neq\X\neq\Sn$, $\Rg=\Sn\setminus\X$;
likewise $\Rg=(\Sn/\Gamma)\setminus\X$ on every spherical space form.
For $|\X|\ge2$ under the spread hypothesis $\dX<\inj$ one also gets
$\Rx=\Sn\setminus\X$ (Corollary~\ref{cor:spread}).
\end{corollary}

\begin{example}[A single point on $\Sn$: the antipodal centre]
\label{ex:point-sphere}
$\X=\{x_{0}\}\subseteq\Sn$. Then $\Mx=\varnothing$, so the two-point
theory reconstructs nothing. But the antipode $\bar x_{0}$ is a medial
centre in $\Mg$: it is joined to $x_{0}$ by an $(n-1)$-sphere of
minimising geodesics, all of length $\pi=\dX(\bar x_{0})$, and the free
ball $B(\bar x_{0},\pi)=\Sn\setminus\{x_{0}\}$ swallows everything.
Corollary~\ref{cor:sphere} thus reconstructs the whole complement from
a single point --- the extreme demonstration that $\Mg$, not $\Mx$, is
the reconstruction axis. The capture curve from any $p\ne x_{0},\bar
x_{0}$ ascends the meridian away from $x_{0}$ and reaches $\bar x_{0}$
in time $\pi-\dX(p)$: capture, class (1).
\end{example}

\begin{example}[Cut-shielding and the two-point axis]\label{ex:shield}
On a prolate ellipsoid of revolution, or on $\Sn$ with a sparse $\X$,
$\Mg$ can strictly contain $\Mx$: a capturing centre $a$ may have two
minimising geodesics to a \emph{single} foot (a cut point of $a$
before any second foot appears), so $a\in\Mg\setminus\Mx$ and the
point it captures is reconstructed by $\Rg$ but not by $\Rx$. This is
the \emph{cut-shielding} obstruction, and it is why the two-point
theorem needs a spread hypothesis (Corollary~\ref{cor:spread}) while
the two-geodesic theorem never does. The obstruction is quantitative:
on prolate ellipsoids it switches on past an eccentricity threshold
(numerically near $c^{\ast}\approx2.6$; see \S\ref{sec:open}).
Class-(1) points of $\Rg$ whose capturing centres lie in
$\Mg\setminus\Mx$ are exactly the shielded points.
\end{example}

\subsection{Hyperbolic space and Hadamard manifolds}
\label{sec:hadamard}

Let $M$ be Hadamard. The ideal medial limit functions are Busemann
functions at a level (Remark~\ref{rem:offset}), so
Theorem~\ref{thm:char} reads: $\Rg$ is the union of the finite medial
balls and the shrunken horoballs $\{b_{\xi}<-c(\xi)\}$ over the
defective ideal points $\xi$. This is the exact Hadamard analogue of
Corollary~\ref{cor:flat}, with ``half-space'' replaced by
``horoball'' and ``$\conv$'' by the horoconvex hull.

The \emph{faithful} question --- which ideal points $\xi$ are
defective, and at what level --- is the curved analogue of
Bia\l{}o\.zyt's defective-hyperplane certificate. In real hyperbolic
space and, more generally, on Hadamard manifolds admitting a faithful
horoconvexity certificate, the answer mirrors the flat one: $\xi$ is
defective iff the contact set $\X\cap\partial\Hb{}$ on the maximal free
horosphere fails to be horoconvex, i.e.\ $\X\cap\partial\Hb{}\neq
\hconv(\X\cap\partial\Hb{})$, and the level is then $0$. We record
this as the natural specialisation; its proof is the horoconvex
counterpart of Motzkin's theorem and is beyond the flat certificate we
assume here (see \S\ref{sec:open} and the discussion in
Remark~\ref{rem:offset}).

\begin{example}[Two ideal directions in $\HH^{2}$]\label{ex:hyp}
Let $\X=\{x_{0},x_{1}\}$ in the hyperbolic plane. The bisector is a
complete geodesic with two ideal endpoints $\xi_{\pm}\in\bdry\HH^{2}$.
Escaping bisector centres $a_{k}\to\xi_{\pm}$ have reaches
$r_{k}=d(a_{k},x_{0})\to\infty$, and $\psi_{a_{k}}=d(\cdot,a_{k})-r_{k}
\to b_{\xi_{\pm}}$, a Busemann function at level $0$; by the swallow
lemma the two open horoballs $\{b_{\xi_{\pm}}<0\}$ are reconstructed.
Together with the finite bisector balls, which fill a lens-shaped
region around the geodesic segment $[x_{0},x_{1}]$, these exhaust
$\Rg$: by Theorem~\ref{thm:char} every reconstructed point lies in a
finite lens ball or in one of the two horoballs. This is the
hyperbolic mirror of Example~\ref{ex:twopoints}, with ``open
half-plane'' replaced by ``open horoball'' --- the escaping centres now
run to the two \emph{ideal points} of the bisector rather than off to
Euclidean infinity, and the swallowed regions are horoballs there. The
contact set on the maximal free horosphere at each $\xi_{\pm}$ is the
single pair $\{x_{0},x_{1}\}$, non-horoconvex, which is the horoconvex
analogue of the defective supporting line and certifies the defect.
Whether the two horoballs and the lens leave any point uncovered ---
the hyperbolic counterpart of the two collinear outer rays --- depends
on the precise horoball geometry and is a faithful-layer question
(\S\ref{sec:open}); the swallow and characterisation above are
unconditional.
\end{example}

\begin{example}[The half-space limit is genuinely ideal]\label{ex:idealneeded}
In both Examples~\ref{ex:twopoints} and \ref{ex:hyp} the drifted
region is reconstructed by \emph{no} finite ball: a point $p$ deep in
the open half-plane / horoball, at an obtuse angle to the bisector,
has its capture curve escaping with a single foot
(Remark~\ref{rem:class2}), so it is class (2), reconstructed only by
the ideal limit. This is why Theorem~\ref{thm:char} must include the
ideal members and why the pure capture theorem
(Corollary~\ref{cor:compact}) cannot hold on a non-compact manifold.
\end{example}

\subsection{Non-compact positive curvature}\label{sec:poscurv}

Here reconstruction is genuinely partial and neither the flat nor the
Hadamard certificate applies; the trichotomy still governs.

\begin{example}[Paraboloid of revolution]\label{ex:paraboloid}
Let $M$ be a paraboloid of revolution ($\sect>0$, a single flattening
end) and $\X$ a small circle about the vertex. The vertex is a finite
medial centre reconstructing the inner disk: class (1). Every point on
the flaring side ascends its meridian forever with a single foot, and
no medial limit function reaches it: class (3). Hence
$\Rg\subsetneq M\setminus\X$ --- partial reconstruction, decided by the
non-compactness of $M$, not by the sign of the curvature.
\end{example}

\begin{example}[Ideal centres in positive curvature]\label{ex:posideal}
Take $\X=\{x_{0},x_{1}\}$ on the paraboloid. The bisector is unbounded,
and its escaping centres furnish genuinely \emph{ideal} medial limit
functions on a $\sect>0$ manifold --- consistent with the trichotomy,
because their sublevel sets are unbounded \emph{non-convex} free
regions. It is the maximal free \emph{convex} regions that are bounded
in positive curvature, not the carriers of medial limit functions; the
distinction matters, and the example shows the ideal class is not a
nonpositive-curvature artefact.
\end{example}

\subsection{The completed picture}\label{sec:table}

Two independent axes organise the theory. \emph{Compactness} decides
partial versus total; the \emph{sign of the Busemann/Hessian data}
decides the flavour of the faithful certificate. The following table
records the two-geodesic theory, with the two-point refinement in
brackets.

\begin{center}
\begin{tabular}{l|ccc}
\toprule
 & $\sect\le0$ & $\sect\equiv0$ & $\sect>0$\\
\midrule
non-compact & partial; $\hconv$ defect & partial; $\conv$ defect
 & partial (paraboloid)\\
 & at level $c(\xi)$ & (Bia\l{}o\.zyt \cite{Bia})
 & finite $+$ ideal centres\\
\midrule
compact & \emph{total} $\Rg$ & \emph{total} $\Rg$
 & \emph{total} $\Rg$\\
 & [$\Mx$: shielding] & [$\Mx$: shielding] & [$\Mx$: shielding;
 $=\Sn$ ok]\\
\bottomrule
\end{tabular}
\end{center}

\noindent The compact row is Corollary~\ref{cor:compact} ---
unconditional, all curvatures, $|\X|\ge1$. The non-compact row is
Theorem~\ref{thm:char}, with the faithful certificate known on the
flat column (Bia\l{}o\.zyt) and its Hadamard extension, and open in
mixed curvature (\S\ref{sec:open}).

\section{Discussion: the two layers}\label{sec:discussion}

It is worth stating plainly what has been achieved and what the shape
of the theory is, because the two are easy to conflate.

\emph{The unconditional layer is complete.} On every complete
Riemannian manifold and for every nonempty closed proper $\X$, the
reconstruction $\Rg$ is exactly the union of the sublevel sets of the
medial limit functions (Theorem~\ref{thm:char}); every point is
captured, drift-swallowed, or escaping (Theorem~\ref{thm:trichotomy});
and reconstruction is total on every compact manifold
(Corollary~\ref{cor:compact}). No curvature bound, genericity,
analyticity, spread, or cardinality hypothesis appears anywhere in
this layer. The two ideas that make it work are worth isolating:

\begin{itemize}
\item \emph{Stop at the axis, do not flow through it.} The margin one
would like to accumulate is not monotone through multi-geodesic points
(Remark~\ref{rem:notmonotone}); but the axis is the \emph{goal}, so
stopping at its first approach both avoids the non-monotonicity and
delivers capture. This is what removes the dependence, in every
previous compact argument, on an unproven extension of the swallow
property past cut points.
\item \emph{Normalise limits by the reach.} Measuring a free ball by
its depth $\dX(a)$ rather than by the distance to a fixed basepoint
makes the limit objects basepoint-free and absorbs the historical
``offset'' into the function itself (Remarks~\ref{rem:horo},
\ref{rem:offset}). The swallow lemma then costs three lines on
\emph{any} complete manifold.
\end{itemize}

\emph{The faithful layer is curvature-specific and only partly known.}
Replacing ``$\psi$ is a medial limit function'' by a certificate
readable from $\X$ and the geometry of $M$ is exactly the content of
Bia\l{}o\.zyt's theorem in the flat case (Corollary~\ref{cor:flat}),
of the horoconvex-defect certificate on Hadamard manifolds
(\S\ref{sec:hadamard}), and it is vacuous on compact manifolds. In
mixed curvature it is open (\S\ref{sec:open}). The honest summary is:
the \emph{dynamics} of reconstruction are understood in full
generality; the \emph{combinatorics} of which asymptotic directions
carry a defect are understood exactly where a Motzkin-type theorem is
available, and Bia\l{}o\.zyt's theorem is the flat instance of such a
theorem.

\section{What remains}\label{sec:open}

The open problems live entirely on the faithful (escape) side.

\begin{enumerate}[label=(\arabic*)]
\item\emph{Certificates in mixed curvature.} For which ideal points
$\xi$ (horofunctions $h$) of a general non-compact $M$ is $h+c$ a
medial limit function, read from the contact structure of $\X$ on
$\{h=0\}$? Known: the $\conv$/Motzkin certificate in $\R^{n}$
(Bia\l{}o\.zyt \cite{Bia}, Motzkin \cite{Mot}); the horoconvex
certificate on Hadamard manifolds. Unknown already for a complete
surface carrying both signs of curvature.
\item\emph{The level function.} On an inhomogeneous Hadamard manifold,
characterise $c(\xi)=\inf\{c:\ b_{\xi}+c\ \text{a medial limit}\}$ from
$\X$ alone (Remark~\ref{rem:offset}). One has $c(\xi)=0$ under
curvature pinching, in the flat case, and over any bisector witness on
the horosphere, and no example with $c(\xi)>0$ is known; deciding
whether the deep-part phenomenon occurs at all is the sharpest open
question on the Hadamard side.
\item\emph{The two-point axis in variable curvature.} Quantify
cut-shielding (Example~\ref{ex:shield}): the prolate-ellipsoid
threshold $c^{\ast}\approx2.6$ awaits an analytic determination, as
does the class of compact $M$ with $\Rx=M\setminus\X$ for all sparse
$\X$ (constant curvature suffices; bounded eccentricity
conjecturally).
\item\emph{Chebyshev-type problems on the frontier.} The Motzkin
question in Carnot groups (e.g.\ $\mathbb{CH}^{n}$ horospheres) and in
Alexandrov spaces of indefinite curvature (pinched horospheres); and
the domain (Lieutier) problem on $\Sn$, where the boundary's own
positively curved intrinsic geometry becomes the detection object.
\end{enumerate}

\paragraph{Conclusion.} Bia\l{}o\.zyt's Theorem~6 generalises to every
complete Riemannian manifold, in two layers. The \emph{unconditional}
layer, proved here in full from that theorem alone, says that
reconstruction by the two-geodesic axis is exactly the union of the
sublevel sets of the medial limit functions
(Theorem~\ref{thm:char}); that every point is captured, swallowed, or
escaping (Theorem~\ref{thm:trichotomy}); and that capture is total on
every compact manifold (Corollary~\ref{cor:compact}), with no
hypothesis beyond completeness and closedness. The \emph{faithful}
layer --- replacing ``is a medial limit function'' by a certificate on
contact sets --- is exact on the flat knife-edge ($\conv$, via
Bia\l{}o\.zyt), on Hadamard manifolds ($\hconv$), and trivial on
compact manifolds; its extension into mixed curvature, and the
Chebyshev-set problems it generates, are the genuine residue. The two
analytic obstructions that historically blocked the Riemannian passage
--- the containment step past cut points and the radial offset --- are
removed here, the first by stopping the ascent at its first approach to
the axis instead of flowing through it, the second by normalising the
limit objects by their own reach.

\begin{thebibliography}{9}

\bibitem{Bia} A.~Bia\l{}o\.zyt, \emph{Sets reconstructible with medial
axis}, preprint, AGH University of Krak\'ow, 2024.

\bibitem{AC} J.-P.~Aubin and A.~Cellina, \emph{Differential
Inclusions}, Grundlehren 264, Springer, 1984. (Existence and extension
for upper semicontinuous convex compact right-hand sides, \S2.1.)

\bibitem{ACa} P.~Albano and P.~Cannarsa, \emph{Structural properties
of singularities of semiconcave functions}, Ann.\ Scuola Norm.\ Sup.\
Pisa \textbf{28} (1999), 719--740.

\bibitem{BH} M.~Bridson and A.~Haefliger, \emph{Metric Spaces of
Non-Positive Curvature}, Grundlehren 319, Springer, 1999.
(Horofunctions and Busemann points, II.8.)

\bibitem{Gro} M.~Gromov, \emph{Hyperbolic manifolds, groups and
actions}, in: Riemann Surfaces and Related Topics, Ann.\ of Math.\
Stud.\ 97, Princeton, 1981, 183--213. (The horofunction
compactification.)

\bibitem{Lie} A.~Lieutier, \emph{Any open bounded subset of $\R^{n}$
has the same homotopy type as its medial axis}, Computer-Aided Design
\textbf{36} (2004), 1029--1046.

\bibitem{MM} C.~Mantegazza and A.~C.~Mennucci, \emph{Hamilton--Jacobi
equations and distance functions on Riemannian manifolds}, Appl.\
Math.\ Optim.\ \textbf{47} (2003), 1--25. (Semiconcavity and first
variation of the distance from a closed set.)

\bibitem{Mot} T.~Motzkin, \emph{Sur quelques propri\'et\'es
caract\'eristiques des ensembles born\'es non convexes}, Atti Accad.\
Naz.\ Lincei \textbf{21} (1935), 773--779.

\end{thebibliography}

\end{document}
